%-------------------------
% Resume in Latex
% Author : Sourabh Bajaj
% Modified for: 林一 (AI Engineer Focus)
% Traditional Chinese Version
%------------------------

\documentclass[letterpaper,11pt]{article}

\usepackage{latexsym}
\usepackage[empty]{fullpage}
\usepackage{titlesec}
\usepackage{marvosym}
\usepackage[usenames,dvipsnames]{color}
\usepackage{verbatim}
\usepackage{enumitem}
\usepackage{fancyhdr}
\usepackage[xetex]{hyperref}
\usepackage{fontspec}
\usepackage{xeCJK}
\usepackage{fontawesome5} % For GitHub icons
% 極限邊距設定：上下左右留白縮小，確保內容能塞進一頁且不擁擠
\usepackage[margin=0.35in]{geometry}

%-------------------------------------------
% Font Settings
% 如果在 Overleaf，請使用 XeLaTeX 編譯器
\setmainfont{Times New Roman}
\setCJKmainfont{Noto Sans CJK TC} 

%-------------------------------------------
% Formatting Setup
\pagestyle{fancy}
\fancyhf{} 
\fancyfoot{}
\renewcommand{\headrulewidth}{0pt}
\renewcommand{\footrulewidth}{0pt}

\urlstyle{same}

\raggedbottom
\raggedright
\setlength{\tabcolsep}{0in}

% 調整 Section 標題下方的間距
\titleformat{\section}{
  \vspace{-8pt}\scshape\raggedright\large
}{}{0em}{}[\color{black}\titlerule \vspace{-5pt}]

%-------------------------
% Custom commands
\newcommand{\resumeItem}[2]{
  \item\small{
    \textbf{#1}{: #2 \vspace{-2pt}}
  }
}

% 無標題的項目 (用於專案細節)
\newcommand{\resumeItemNoTitle}[1]{
  \item\small{
    {#1 \vspace{-2pt}}
  }
}

% 調整 Subheading 表格寬度與間距
\newcommand{\resumeSubheading}[4]{
  \vspace{-1pt}\item
    \begin{tabular*}{0.97\textwidth}{l@{\extracolsep{\fill}}r}
      \textbf{#1} & #2 \\
      \textit{\small#3} & \textit{\small #4} \\
    \end{tabular*}\vspace{-5pt}
}

\newcommand{\resumeSubHeadingListStart}{\begin{itemize}[leftmargin=*]}
\newcommand{\resumeSubHeadingListEnd}{\end{itemize}}
\newcommand{\resumeItemListStart}{\begin{itemize}}
\newcommand{\resumeItemListEnd}{\end{itemize}\vspace{-5pt}}

%-------------------------------------------
%%%%%%  CV STARTS HERE  %%%%%%%%%%%%%%%%%%%%%%%%%%%%

\begin{document}

%----------HEADING-----------------
% 聯絡資訊區塊
\begin{tabular*}{\textwidth}{l@{\extracolsep{\fill}}r}
  \textbf{\href{https://github.com/tana0101}{\Large 林 一}} & Email: \href{mailto:xenone0101@gmail.com}{xenone0101@gmail.com}\\
  \href{https://github.com/tana0101}{https://github.com/tana0101} & Mobile: (+886) 9-7132-9517 \\
\end{tabular*}

%-----------SUMMARY-----------------
\section{個人簡介 (Summary)}
目前就讀於成大資工所碩士，專注於電腦視覺與深度學習。具備將軟體工程原則導入 AI 開發的實務經驗，曾獲 \textbf{2025 台積電 IT CareerHack 冠軍}，並以\textbf{第一作者}於 \textbf{IEEE ICMAIL 2023} 發表論文。

%-----------EDUCATION-----------------
\section{教育背景 (Education)}
  \resumeSubHeadingListStart
    \resumeSubheading
      {國立成功大學 (NCKU)}{台南，台灣}
      {資訊工程學系 碩士班}{2024 年 9 月 - 預計 2026 年 7 月}
      \resumeItemListStart
        \resumeItem{碩士論文}
          {基於深度學習之小兒髖關節 X 光分析系統}
        \resumeItem{相關修課}
          {深度學習、電腦視覺、影像處理}
      \resumeItemListEnd
      
    \resumeSubheading
      {國立臺灣海洋大學 (NTOU)}{基隆，台灣}
      {資訊工程學系 學士班 (GPA: 3.77/4.0)}{2020 年 9 月 - 2024 年 6 月}
  \resumeSubHeadingListEnd

%-----------SELECTED PROJECTS-----------------
\section{精選專案 (Selected Projects)}
  \resumeSubHeadingListStart

    % Project 1: Master Thesis (SOTA & Clinical Impact focus)
      \resumeSubheading
      {基於深度學習之小兒髖關節 X 光分析系統 | 關鍵點偵測}{碩士論文}
      {PyTorch, YOLO, ConvNeXt, SimCC, Medical Imaging}{\href{https://github.com/tana0101/Hip-Joint-Keypoint-Detection}{\faGithub \textbf{ GitHub 程式碼}}}
      \resumeItemListStart
        \resumeItemNoTitle
          {\textbf{建構}髖關節發育不良 (DDH) 的診斷流程，採用\textbf{二階段由上而下 (top-down)}的關鍵點回歸方法；整合 \textbf{YOLO} 進行 ROI 定位，並採用 \textbf{ConvNeXt 結合 SimCC} 實現次像素級的關鍵點回歸。}
        \resumeItemNoTitle
          {在 MTDDH 資料集達成 \textbf{SOTA} 表現，將 \textbf{YOLOv11-seg}（F1: 0.593）提升 \textbf{59\% 至 0.943}，以提高臨床的診斷準確率。}
        \resumeItemNoTitle
          {經臨床驗證，系統誤差低於資深醫師 (\textbf{平均誤差：5.01px vs. 6.34px})，顯示其在\textbf{自動化診斷}上的精準度與臨床實用價值。}
        \resumeItemNoTitle
          {設計可擴充的\textbf{模組化架構}，支援多種 backbone（HRNet、EfficientNet）與 head 設計，用於髖臼指數（AI）角度量測與 IHDI 分類。}
      \resumeItemListEnd

    % Project 2: TSMC Hackathon (First Place focus)
    \resumeSubheading
      {企業溝通無國界翻譯系統 (2025 台積電 IT CareerHack)}{第一名 (冠軍)}
      {JavaScript, String Matching Algorithm, Frontend Development}{\href{https://github.com/daidaiclub/Careerhack-2025-RealTime-Translate-System}{\faGithub \textbf{ GitHub 程式碼}}}
      \resumeItemListStart
        \resumeItemNoTitle
          {開發企業用即時翻譯工具，可\textbf{自動標示}特定領域的關鍵術語；前端採用 \textbf{HTML/CSS/JS} 製作，並榮獲\textbf{冠軍}。}
        \resumeItemNoTitle
          {設計並實作\textbf{術語比對演算法}，將翻譯內容中的技術名詞\textbf{自動對應與註記}，提升閱讀清楚度與資訊一致性。}
        \resumeItemNoTitle
          {向台積電主管進行成果展示與簡報，展現良好的技術溝通與表達能力。}
      \resumeItemListEnd

    % 3. Undergraduate Research (NTOU)
    \resumeSubheading
      {基於深度學習之水下影像還原系統}{大學專題研究}
      {PyTorch, GAN, Image Enhancement}
      {
        \href{https://github.com/NTOU-Arrays-Start-at-One/Give-ocean-a-piece-of-your-mind}{\faGithub \textbf{ 程式碼}} 
        \quad % 加一點間距
        \href{https://doi.org/10.1109/ICMAIL59311.2023.10347502}{\faFilePdf \textbf{ 論文}}
      }
      \resumeItemListStart
        \resumeItemNoTitle
          {以\textbf{第一作者}身份於 \textbf{IEEE ICMAIL 2023} 發表研討會論文。}
        \resumeItemNoTitle
          {整合 \textbf{WaterNet（CNN 水下顏色恢復模型）}與 \textbf{GAN} 上色模型，建立水下影像色彩還原系統以改善紅光衰減造成的色偏；並使用色板分析與 \textbf{K-means + CIEDE2000} 進行量化評估，驗證色彩還原效果。}
      \resumeItemListEnd

  \resumeSubHeadingListEnd

%-----------EXPERIENCE-----------------
\section{實務經驗 (Experience)}
  \resumeSubHeadingListStart

    \resumeSubheading
      {Programing.tw}{台灣}
      {共同創辦人暨產品負責人}{2021 年 2 月 - 至今}
      \resumeItemListStart
        \resumeItem{全端開發}
          {使用 \textbf{Go (Gin)} 與 \textbf{Vue.js} 建構「Proglearn」平台，並實作 \textbf{WebSocket} 用於即時遠端教學。}
        \resumeItem{獲獎肯定}
          {憑藉創新概念於潮創客創業大賽獲得\textbf{優選獎（前四名）}，並在 2023 資訊智慧創新跨域專題競賽中獲得\textbf{特優獎}。}
      \resumeItemListEnd

  \resumeSubHeadingListEnd

%-----------SKILLS-----------------
\section{專業技能 (Technical Skills)}
 \resumeSubHeadingListStart
    \resumeItem{AI 與電腦視覺}
      {PyTorch, OpenCV, YOLO, ConvNeXt, Data Augmentation, GANs}
    \resumeItem{程式語言}
      {Python, C/C++, Java, Go, JavaScript, SQL}
    \resumeItem{網頁與工具}
      {Vue.js, Gin (Go), Docker, Git, Linux (Ubuntu), LaTeX, CI/CD}
 \resumeSubHeadingListEnd

%-----------AWARDS-----------------
\section{榮譽與獎項 (Honors \& Awards)}
  \resumeSubHeadingListStart
    \resumeItem{程式競賽}
      {2022 ICPC Asia Taiwan Online Programming Contest - 佳作 (前 35\%)}
  \resumeSubHeadingListEnd

\end{document}